\chapter{Piano development studies}\label{app:phase1dev}
This appendix collects the development studies of the piano setting in
full. Each of them has a short summary in the main line
(\S\ref{sec:masksweep}, \S\ref{sec:devstudies}, \S\ref{sec:musictheory}),
and all of the evidence reported here comes from the development set under
the protocol of \S\ref{sec:hygiene}. The studies concern, in turn, the
sensitivity of the results to the masking level, the sufficiency of
hand-built score features, calibrated deep baselines, task-matched
pretraining of the music model, the comparison of graph kernels in the
two-stage regime, and music-theoretic feature columns together with probes
of the embeddings.

\section{The masking-level sweep, in full}\label{app:masksweep}
The published operating point hides \maskRate{} of the notes of each piece.
That fraction was inherited from earlier work rather than derived, so its
influence was measured on the development pieces after the held-out
evaluation. The hidden fraction was swept from 50\% down to 10\%, with the
leak-free mask-aware embeddings recomputed at each fraction, since
leak-freeness requires blanking exactly the notes that each mask hides. The
sweep was extended to the leave-one-out limit, where every note is predicted
from all the others through the closed-form GP identity. At that limit only,
two relaxations of the protocol apply. The hyperparameters are fitted on the
full piece, which is the textbook leave-one-out cross-validation estimator,
and the embeddings are the leak-free pre-velocity readout rather than
per-mask re-embeddings. Both relaxations apply identically to all three
systems, so the comparison of their ordering remains a fair paired
comparison, although the absolute numbers of the leave-one-out row are not
strict-protocol values.

The ordering of the models, with the proposed model below the no-graph
ablation and that below the features-only model, is stable at every level
including leave-one-out, so the operating point is arbitrary but not
fortunate. The paired contribution of the graph is significant from 50\%
down to 20\% hidden and at the leave-one-out limit. The contribution of the
embeddings grows monotonically with observation density, from $-0.005$ RMSE
at 50\% hidden to $-0.021$ at leave-one-out, paired, as context-dependent
features would be expected to behave. Coverage stays at 0.92--0.93
throughout, so calibration is insensitive to the masking level. One
piece-seed combination at 30\% hidden reproduces, on a development piece,
the Gaussian-tail failure of the log score described in
\S\ref{sec:guard-gp}, identically across all three systems, which is
consistent with the failure being a property of the tail of the likelihood
rather than of any one model. Figure~\ref{fig:masksweep} shows both curves.

\figph{masksweep_dev.png}{1.0}{Sensitivity to the masking level on the
development set. Left: pooled held-out RMSE as the hidden fraction falls
from 50\% to 10\% and to the leave-one-out limit (open markers, whose
hyperparameters are fitted on the full piece). Right: the paired per-piece
contributions of the graph and of the music-model embeddings at each level,
with bootstrap 95\% intervals. The interval for the contribution of the
graph excludes zero at every level except 10\% hidden, where only
$\sim$10\% of the notes are scored per piece.\label{fig:masksweep}}

\section{Hand-built score features and the music model}
The second study asks whether hand-built score features suffice in place of
the music model, on the development set. Against the 25-feature hand-built
score baseline fitted under the same protocol, the music model shows no
measurable difference in average error. Its edge, which is significant,
lies in calibration and in the velocity channel. The two sources of
information also combine. Under the two-stage plug-in baseline,
concatenating the two means lowers the RMSE from $\rmseNetGraph$ to
$\rmseStack$, and in the proposed model, where both enter as feature
kernels, the marginal value of the embeddings is $-0.008$ RMSE and $-0.034$
NLL, paired and significant, with the cost of their removal reproduced in
direction on the held-out test set. Per channel, the mean-only edge of the
music model over the 25-feature baseline is confined to velocity, with no
measurable difference on timing and articulation. This is the same
placement that the decomposition of \S\ref{sec:attribution} recovers from
inside the fitted model.

\section{Calibrated deep baselines}
The comparisons of the main line give the deep model a single shared
predictive scale, which understates it as a rival on the calibration axis.
Two deep baselines that produce their own per-note uncertainty were
therefore evaluated on the development set under the identical strict
protocol and information set: the 25 features together with the music-model
embeddings, standardized per piece, with strict mask-aware embeddings at
evaluation and the same per-piece variance floor. The first is a
heteroscedastic MLP head, which outputs $\mu_c(x)$ and $\log\sigma_c^2(x)$
per channel and is trained by Gaussian negative log-likelihood on the head
pieces. The second is a five-member deep ensemble of such heads. The
proposed model is better than both on both axes by wide paired margins. Its
RMSE is $0.360$ against $0.450$--$0.454$, a paired difference of $-0.090$
with interval $[-0.115, -0.067]$, and its NLL is $-0.404$ against $-0.069$
to $-0.104$, a paired difference of $-0.30$ with interval $[-0.40, -0.21]$.
Ensembling five heads adds almost nothing over one.

Where the deep heads land is the instructive part. At an RMSE of
$\approx 0.45$ they sit at the level of the cross-piece plug-in mean
without a graph. Their coverage is respectable at $0.92$, so the variance
head does its job, but their NLL is far behind, because ranking which notes
are uncertain requires per-piece adaptation and the graph. This is the
attribution result of the main line, reached from the other side. A
follow-up at the leave-one-out limit of \S\ref{sec:masksweep}, the densest
observation and the regime in which a read-out head should be at its best,
leaves the conclusion unchanged. The proposed model is better on every
channel on both axes: against the ensemble, timing RMSE is $0.149$ against
$0.165$ and NLL $-1.12$ against $-0.50$, articulation RMSE is $0.546$
against $0.755$, and velocity RMSE is $0.066$ against $0.084$. The
uncertainty ranking of the ensemble therefore does not catch up with denser
observation. The scope of the comparison is that of read-out heads on fixed
embeddings, the natural deep comparator at matched information and training
data, and a fine-tuned end-to-end deep model remains untested.

\section{Task-matched pretraining}
A bidirectional, mask-aware retraining of the music model, matched exactly
to the evaluation task, was compared with the original at matched compute on
the development set. It is not better. Pooled error is slightly but
significantly worse, by $+0.011$, confidence quality shows no measurable
difference strict-to-strict ($-0.325$ against $-0.322$), and a $3\times$
training budget does not close the gap, since timing improves while
articulation and velocity do not. The reading is that the original read-out
had already extracted what the change of objective was meant to unlock. The
study also showed that the bidirectional read-out admits the leakage pattern
that the causal read-out excludes, because a note read at its own velocity
token contains that velocity. The pattern was recognised by its signature,
and the design rule that follows from it is a leave-one-out read-out, in
which each note is read with its own tokens masked, enforced by a
structural test.

\section{Kernel comparison in the two-stage regime}\label{sec:kernels}
The kernel comparison was run under the two-stage plug-in baseline on the
development set: the plug-in mean $\muLM$ under the strict protocol, one
scalar GP per channel, identical masks in every row, and recovery and
calibration scored paired per piece. It moved one of the two design axes of
\S\ref{sec:graphvariants} at a time, either the graph $\Lg$ or the spectral
profile $g(\nu;s)$. Three findings stand, and they are the reason the
proposed model keeps the plain graph with the additive profile.

First, the profile is inert. Graph Mat\'ern at $\alpha = 2$ and $3$, the
diffusion kernel and the normalized Laplacian are statistically
indistinguishable from the additive Laplacian on RMSE ($0.392$--$0.393$
against $0.393$) and do not improve NLL. Their evidence-fitted filters
coincide over the low-frequency components, where the prior mass lies, and
differ only where the feature kernels and the noise absorb the variance.
Second, the graph is not inert. Removing every edge raises RMSE to $0.446$
and NLL to $-0.190$, both paired differences with intervals excluding zero,
and the temporal chain shows no measurable difference from the full graph
on RMSE but is worse calibrated, with a paired NLL difference of $+0.031$
whose interval excludes zero. Third, music theory helps as extra edges and
not as a replacement metric. Tonal distance in place of semitone distance
is worse on both axes (RMSE $0.403$, NLL $-0.300$, both intervals excluding
zero), whereas adding chord and voice-leading edges gave the only
improvement on both axes, a paired RMSE difference of $-0.008$ and a paired
NLL difference of $-0.013$ with intervals excluding zero, robust to a sweep
over six settings of the untuned edge weights scaled by $\times 3$ and by
$\div 3$.

Under the proposed model the chord edges are redundant at the 40\%
operating point once the music-model embeddings enter the kernel. The
effect follows a density gradient, with a recovery gain at 30/20/10\%
hidden, no measurable difference at 40\%, and nothing at 50\%. The proposed
model therefore keeps the plain graph, and adopting the chord edges would
call for a fresh held-out evaluation under the protocol of
\S\ref{sec:hygiene}. The full tables, the row definitions and the $\mu = 0$
block are in the development record.\footnote{The kernel comparison is
recorded in \texttt{docs/kernel\_comparison\_results.md} and the density
gradient of the chord edges in \texttt{docs/kernel\_multirate\_results.md}.}

\section{Music-theoretic feature columns and embedding probes}\label{app:theoryfeat}
The incremental route adds music-theoretic score features to the hand-built
set, where they compose with any kernel. Fourteen score-only columns were
added: local key by Krumhansl--Schmuckler profiles (key clarity, mode, a
scale-degree encoding, an in-scale flag and the circle-of-fifths motion of
the local key), a metrical-weight hierarchy, within-chord dissonance and a
bass flag, phrase-boundary salience, motif-repetition counts and per-voice
tessitura. On the development set they change nothing of note. The paired
RMSE difference is $-0.002$ without the music-model embeddings, with an
interval that touches zero, and $+0.001$ beside them, with an interval
covering zero, and no single channel is significant in either regime. The
finding is not specific to the 40\% operating point. Repeated at every
masking level of \S\ref{sec:masksweep}, the columns stay non-significant
beside the embeddings at each level and reach marginal significance on the
plain model only at 10\% hidden, with a paired RMSE difference of $-0.004$.
This is the same pattern as the chord edges show: where the per-piece
weights are best determined, explicit theory carries a small signal that
the embeddings already cover.

A follow-up linear-probe study sharpens the interpretation, and not in the
obvious direction. Probing the embeddings directly, with a cross-piece
ridge fitted on the head pieces and scored out of sample on the development
pieces, shows that they encode voicing and rhythm: bass-note membership
(AUC $0.99$), metrical weight, phrase-boundary salience and register. They
do not encode tonality. Scale degree, in-scale membership, mode and local
key are all at chance from the embeddings ($R^2 \approx 0$, AUC
$\approx 0.5$), whether or not piece-level information is removed, and a
2048-dimensional random-features probe replicates every tonal result, so
this is not an artifact of probe linearity. The failure of the theory
columns is therefore \emph{not} explained by the embeddings already
carrying the tonal signal. Taken together with the tonal-metric result of
route (A) of \S\ref{sec:musictheory}, the consistent reading is that tonal
structure carries no measurable marginal signal for these three expressive
channels at this scale, while the rhythmic and voicing structure that the
columns add is already in the kernel. This sharpens the attribution of
\S\ref{sec:attribution} in an unexpected direction: the value of the
embeddings lies in rhythm, texture and dynamics convention, not in harmony.
The columns remain available in the implementation as an interpretable
probe and are switched off by default.

% =====================================================================
